\subsection{Important functions}
\label{subsec:important_functions}
\subsubsection{Distance calculation: Haversine distance}
This is the formula to compute user-to-venue distance from latitude and longitude.
\subsubsection{UI: Tailwind CSS, ShadCN}
To improve the UI and speed up development, we adopted Tailwind CSS and the
ShadCN UI components \cite{shadcn_ui}. They are used throughout the frontend. We also modified
these components in order to match our desired functionalities, such as the data tables, which
we used in many of our pages. On top of this, we used tweakCN's Clean Slate theme \cite{tweakcn} and
shadcn.io's component \cite{shadcn_io} for dark mode transition animation.
\subsubsection{Useful React hooks}
Here are some React hooks we designed that were very useful:
\begin{itemize}%[leftmargin=*]
    \item \verb|use-async|: This is used to determine whether the foreground and background should be loading.
    Foreground loading is to show the loading animation (once \(> 300\) ms), background loading
    is to prevent user spamming HTTP requests. See section \ref{subsubsec:ux_loading_screen} and 
    \ref{subsubsec:security_prevent_rapid_request} for more information. This also
    gives us a systematic way to determine the last sync time.
    \item \verb|use-message|: This is a useful hook for showing a message and determining its color.
    See section \ref{subsubsec:ux_feedback_messages} for more information.
\end{itemize}
\subsubsection{Centralized API Gateway}
Our frontend was very messy when it comes to backend API requests. As we have \verb|fetch()| everywhere
and each with a slightly different implementation, we see a lot of small bugs and crashes due to improper error handling.
To solve this, we finally wrote a centralized function for all backend requests:
\begin{verbatim}
    export async function requestToBackend(method, endPoint, jsonData = null)
\end{verbatim}

